\documentclass[11pt]{letter}
\usepackage[utf8]{inputenc}
\usepackage{geometry}
\geometry{letterpaper, margin=1in}

\address{Rafael D. De Paz \\ Independent Researcher \\ me@rdepaz.com}
\signature{Rafael D. De Paz}

\begin{document}
\begin{letter}{Editorial Board \\ Physical Review D \\ American Physical Society}

\opening{Dear Editors,}

I submit for your consideration the manuscript \textit{"The Cosmological Parsimony Exploit: Resolving the Hubble Tension and the Deterministic Evolution of the Minimal Acceleration Scale ($a_0$)"}.

In this paper, I formalize a solution to the "Cosmological Coincidence Problem"—the unexplained numerical correlation between the local galactic acceleration threshold ($a_0$) and the universal expansion rate ($H_0$). Mainstream cosmology currently assumes this $10^{-60}$ scale alignment is an accident, allowing the Dark Matter hypothesis to persist. By enforcing the strict mathematics of Algorithmic Parsimony (The Logos), this paper establishes that "coincidences" of this magnitude are structurally impossible in a deterministic framework. 

Consequently, we mathematically derive that $a_0$ is mechanically bound to $H(z)$ and must dynamically evolve over cosmic time. This model explicitly removes the theoretical necessity (and compounding theoretical debt) of the Cold Dark Matter particle, directly predicting specific high-redshift galactic rotational anomalies resolvable by upcoming James Webb Space Telescope (JWST) surveys.

This research rigorously bridges theoretical astrophysics, algorithmic information theory, and the Mode Identity Theory to halt the current epistemological drift in the $\Lambda$CDM standard model. I believe it is a vital theoretical addition to your journal.

\closing{Sincerely,}

\end{letter}
\end{document}
